\chapter{Discussion}
\section{Overall Findings}
\par Based on the above experimental results, this thesis systematically
compares the performance of classical machine learning methods and a
lightweight convolutional neural network in the patch-based image focus
detection task. Overall, the experiments show that under a unified data
pipeline, unified automatic labeling rules, and unified evaluation criteria,
different methods indeed present clear and stable performance differences.
However, the more important conclusion is not that one class of models achieves
the highest metric in a single experiment, but that under the current
experimental setting of this thesis, the performance gain brought by increasing
the training sample ratio from $20\%$ to $100\%$ is generally limited. Under
both the binary and three-class settings, the metric changes of most models are
relatively moderate, and are reflected more as slight fluctuation or limited
improvement, while the training time increases clearly. It should be emphasized
that the ``sample size'' here mainly refers to the ratio of patch training
samples obtained by splitting a limited number of original images, rather than
the unlimited expansion of independently collected image quantity. Therefore,
the result of this thesis is better understood as follows: under the current
conditions of $32\times32$ patches, automatic labeling rules, and fixed data
source, continuing to enlarge the training sample ratio changes the overall
performance pattern only to a limited extent.

\par From the perspective of model performance trends, this thesis observes
that the advantages of different methods are closely related to task
definition. In the binary task, although the CNN achieves the highest overall
sharp-class metric, the SVM and random forest can also maintain stable
performance at relatively close evaluation levels, and their training cost is
lower. In the three-class task, the results of random forest and SVM further
show that traditional methods also have strong applicability in the current
task. By combining the complete sample-size scan results, it can be seen that
model performance is related not only to structural complexity, but also to
the boundary characteristics of the task itself. In the setting of this
thesis, patch labels are generated according to local sharpness-related
information, and both traditional models and CNN are trained on unified
grayscale patch input. The experimental results show that this patch
representation is already sufficient to support classical machine learning
methods in learning competitive decision boundaries. Under the binary setting,
the boundary between sharp and non-sharp is relatively direct. Under the three-
class setting, however, the boundary between the sharp class and the
intermediate-state class is clearly more complex. Therefore, under the task
conditions defined in this thesis, machine learning (ML) methods already have
high applicability, but this judgment should be understood as a conclusion
under the current input representation and experimental scope, rather than a
universal summary for all focus detection tasks.

\section{Practical Significance}
\par One important significance of this research is that it provides reference
for model selection in practical applications. For scenarios with limited data
scale, constrained computational resources, or a higher requirement for model
interpretability, classical machine learning methods may still be more feasible
choices \cite{murphy2022pml}, because they usually train faster, require lower
implementation cost, and can maintain relatively stable discrimination effect
under a unified patch input representation. In particular, from the
experimental results, the SVM maintains high recall at relatively low training
cost in the binary task, and the random forest achieves a relatively strong
sharp-class F1-score in the three-class task, while their training times are
both significantly lower than that of the CNN. This indicates that when the
difference in evaluation results is not large, traditional methods often have
higher practical cost-effectiveness in engineering applications, rather than
serving only as baseline models.

\par At the same time, if the application scenario has more available samples,
allows higher training cost, or places more emphasis on the model's automatic
learning ability for complex local patterns, the lightweight CNN
\cite{lecun2002gradient,krizhevsky2012imagenet} still has room for further
optimization. The experiments in this thesis show that the CNN can indeed
obtain relatively strong overall recognition results in the binary task, but
this advantage is accompanied by higher training cost. It should be noted that
the epoch setting of the CNN in the main text is itself the result of a trade-
off between performance and time cost. Its purpose is to ensure that different
experiments can be completed under a unified and acceptable computational cost,
rather than waiting until the model reaches the fullest possible convergence
under every setting. Even so, the CNN still shows relatively high accuracy and
strong overall sharp-class discrimination ability, which indicates that its
automatic representation ability has already shown some advantage in the
current task. However, from the perspective of training time, its cost is still
higher than that of classical machine learning methods. It should also be added
that in the three-class task, the CNN under the default main experimental
configuration performs worse than random forest and SVM, but the supplementary
experiment already shows that training strategies such as weighted sampling can
clearly improve its sharp-class recall and F1-score. Therefore, from a
practical perspective, when the effects of ML and DL are close, whether it is
necessary to invest more computational resources for limited performance gain
still needs to be judged according to the specific application requirement. At
the same time, the evaluation of CNN should not be reduced simply to the
advantages or disadvantages of the model category itself, but should be
analyzed together with the training configuration.

\section{Limitations of the Study}
\par Although this thesis compares multiple kinds of models under a unified
experimental framework, the study still has several limitations. First, the
experimental data scale is relatively limited, and the patch labels are
generated by automatic rules rather than precise manual annotation. Therefore,
the labels themselves may contain some noise. This noise may affect model
training quality and, to some extent, limit the external generalization of the
experimental conclusions. If some models in the experiments show bias in
recall, precision, or class boundary behavior, the reason may partly come from
the approximate partition of the sharp-blur boundary by the automatic labeling
mechanism itself, rather than being determined completely by model capability.

\par Second, the model range adopted in this thesis mainly focuses on decision
tree, random forest, SVM based on Nystr\"{o}m kernel approximation, and a
lightweight CNN, and does not yet include more complex deep network structures
or more advanced representation learning strategies. Therefore, the current
conclusion is more suitable for the low-sample and lightweight comparison
scenario defined in this thesis, and cannot be directly generalized to all
image focus detection tasks. More specifically, the conclusion of this thesis
mainly targets cases where local boundary features are relatively clear and the
patch representation can effectively preserve discriminative information. If
the task is further extended to more complex blur types, stronger semantic
interference, or more difficult cross-domain conditions, the potential
advantage of DL methods may still appear more clearly. In addition, although
the preprocessing pipeline, patch size, and evaluation method are uniformly
controlled in the experiments, the results may still be affected by threshold
setting, class distribution, automatic labeling parameters, and differences in
data source. Therefore, the conclusion of this thesis should be understood as
follows: under the experimental conditions based on $32\times32$ patches,
rule-based automatic labeling, and limited-sample training, ML can already
handle this kind of problem well at relatively low cost, rather than providing
an absolute conclusion that is universally valid for all scenarios.

\section{Future Work}
\par In future work, this research can be further improved from three aspects:
data, model, and task extension. First, at the data level, larger-scale and
more diverse image datasets can be introduced, together with higher-quality
manual annotation or semi-automatic correction mechanisms, so as to improve
the reliability of patch labels and the robustness of experimental conclusions.
If conditions allow, the influence of different data sources, different blur
types, and different patch-size settings on model performance can also be
further compared, so as to verify the applicability of the conclusions of this
thesis in broader scenarios.

\par Second, at the model level, more complex deep network structures than the
lightweight CNN can be explored, or transfer learning, attention mechanisms,
self-supervised learning, and more systematic threshold calibration strategies
can be introduced, so as to examine whether DL methods can bring more
significant gains when task complexity continues to increase. At the same time,
feature design, feature fusion strategy, and class-imbalance handling in
traditional methods can also be further optimized to examine their performance
upper bound under low-cost conditions. Finally, at the task level, future
research can extend patch-level classification to more continuous sharpness
estimation, focus saliency map prediction, or finer focus-region segmentation
tasks, so as to improve the applicability of the method in practical image
quality analysis and visual perception applications. Based on the results of
this thesis, future research can focus on how to further improve the detection
ability of minority sharp regions while maintaining lower labeling cost and
higher training efficiency, and evaluate whether increasing model complexity
and computational resource investment is really necessary in more complex
tasks.
